\documentclass[11pt,a4paper]{article}

% ----------------------------------------------------
% PACKAGES ESSENTIELS
% ----------------------------------------------------
% Nota: inputenc se omite a menudo en compiladores modernos (LuaLaTeX/XeLaTeX),
% pero se mantiene aquí por compatibilidad con el código original si usas PDFLaTeX.
%\usepackage[utf8]{inputenc}     % Encodage du texte (UTF-8)
\usepackage[T1]{fontenc}        % Encodage des caractères
\usepackage[french]{babel}      % Langue du document
\usepackage{geometry}           % Gestion des marges
\usepackage{graphicx}           % Insertion d’images
\usepackage{amsmath, amssymb}   % Symboles et équations mathématiques
\usepackage{siunitx}            % Unités du Système International (SI)
\usepackage{caption}            % Amélioration des légendes
\usepackage{subcaption}         % Sous-figures
\usepackage{booktabs}           % Tableaux esthétiques
\usepackage{float}              % Contrôle de la position des figures
\usepackage{xcolor}             % Gestion des couleurs
\usepackage{fancyhdr}           % En-têtes et pieds de page personnalisés
\usepackage{enumitem}           % Listes personnalisées
\usepackage{physics}            % Notation physique (contraintes, etc.)
\usepackage{titlesec}           % Style des sections
\usepackage[strict]{changepage} % Ajustement des marges locales
\usepackage{framed}             % Cadres et encadrés

% Hyperref doit être chargé en dernier (généralement)
\usepackage{hyperref}           % Liens cliquables dans le PDF

% ----------------------------------------------------
% PARAMÈTRES DE MISE EN PAGE
% ----------------------------------------------------
\geometry{margin=2.5cm}            % Marges du document
\setlength{\parskip}{0.5em}        % Espacement entre paragraphes
\setlength{\parindent}{0pt}        % Pas d’indentation au début des paragraphes

% ----------------------------------------------------
% CONFIGURATION DES LIENS
% ----------------------------------------------------
\hypersetup{
    colorlinks=true,
    linkcolor=enstaBleuFonce, % Utilisation de la couleur ENSTA pour les liens internes (TOC)
    urlcolor=blue,
    citecolor=gray
}

% ----------------------------------------------------
% COULEURS ENSTA
% ----------------------------------------------------
\definecolor{enstaBleuFonce}{HTML}{003366}    % Bleu marine (sections)
\definecolor{enstaBleuClair}{HTML}{0073CF}    % Bleu moyen (sous-sections)
\definecolor{formalshade}{rgb}{0.95,0.95,1}   % Fond clair pour encadrés

% ----------------------------------------------------
% STYLE DES SECTIONS
% ----------------------------------------------------
% Section (bleu marine, majuscule, gras, ligne horizontale)
\titleformat{\section}[block]
  {\normalfont\Large\bfseries\color{enstaBleuFonce}}
  {\thesection}{1em}{}
  [\vspace{0.3em}\titlerule\color{enstaBleuFonce}\vspace{0.3em}]

% Sous-section (bleu clair)
\titleformat{\subsection}
  {\normalfont\large\bfseries\color{enstaBleuClair}}
  {\thesubsection}{1em}{}

% Sous-sous-section (gris)
\titleformat{\subsubsection}
  {\normalfont\normalsize\bfseries\color{black!70}}
  {\thesubsubsection}{1em}{}

% ----------------------------------------------------
% EN-TÊTES ET PIEDS DE PAGE
% ----------------------------------------------------
\pagestyle{fancy}
\fancyhf{}
\fancyhead[L]{ENSTA Paris}
\fancyhead[R]{TP : Comportement non linéaire des matériaux}
\fancyfoot[C]{\thepage}

% ----------------------------------------------------
% ENVIRONNEMENT FORMAL (ENCADRÉ)
% ----------------------------------------------------
\newenvironment{formal}{%
\def\FrameCommand{%
  \hspace{1pt}%
  {\color{enstaBleuFonce}\vrule width 2pt}%
  {\color{formalshade}\vrule width 4pt}%
  \colorbox{formalshade}%
}%
\MakeFramed{\advance\hsize-\width\FrameRestore}%
\noindent\hspace{-4.55pt}%
\begin{adjustwidth}{}{7pt}%
\vspace{4pt}%
}{%
\vspace{4pt}\end{adjustwidth}\endMakeFramed%
}

% ----------------------------------------------------
% DÉBUT DU DOCUMENT
% ----------------------------------------------------
\begin{document}

% ====================================================
% PAGE DE COUVERTURE
% ====================================================
\begin{titlepage}
    \centering
    \vspace*{4cm}
    
    % Remplacer ce bloc par \includegraphics[width=0.6\textwidth]{imgs/logo_ensta_2025.jpg}
    \framebox{\parbox{0.6\textwidth}{\centering
        \vspace{1cm}
        \textbf{LOGO ENSTA} \\
        (Image non trouvée)
        \vspace{1cm}
    }}\par\vspace{0.5cm}
    
    \vspace{0.5cm}
    {\Large Comportement non linéaire des matériaux \par}
    \vspace{0.2cm}
    {\huge\bfseries Travail Pratique \par}
    \vspace{3cm}
    {\Large Sofia PINILLA SALAS \par}
    {\Large Yasmine RABIA \par}
    {\Large  \par}
    \vspace{1.5cm}
    \textbf{Encadrant :} Fabien SZMYTKA \par
    \vfill
    École Nationale des techniques avancées\\
    Majeure Mecanique\\
    Mars 2026 \par
\end{titlepage}

% ====================================================
% TABLE DES MATIÈRES (ÍNDICE)
% ====================================================
\newpage
% Change le nom "Table des matières" si nécessaire avec :
% \renewcommand{\contentsname}{Sommaire}
\tableofcontents
\newpage

% ====================================================
% INTRODUCTION
% ====================================================
\section{Introduction}

Ce premier exercice du travail dirigé a pour objectif de se familiariser avec le logiciel fourni et d’apprendre à réaliser des simulations par éléments finis sur un seul élément pour des essais de traction, de traction-compression et de relaxation. 

\begin{figure}[H]
    \centering
    % Remplacer ce bloc par \includegraphics[width=0.75\linewidth]{imgs/Graph2-1Q5.png}
    \framebox{\parbox{0.75\linewidth}{\centering
        \vspace{2cm}
        \textbf{Graphique 2-1Q5} \\
        (Image Placeholder)
        \vspace{2cm}
    }}
    \caption{Écrouissage }
    \label{fig:ecrouissage}
\end{figure}

Dans cette configuration, les valeurs de $\mathbf{R_0 = 145}$, $\mathbf{Q = 240}$ et $\mathbf{b = 480}$ ont permis d'atteindre la meilleure adéquation pour les phases linéaire et plastique. \\

\begin{table}[H]
    \centering
    \begin{tabular}{l c}
        \toprule
        \textbf{Propriété} & \textbf{Valeur} \\
        \midrule
        Module  ($E$) & \num{158000} \\ 
        Coefficient ($\nu$) & 0.3 \\
        $R_0$ & 145 \\ 
        $Q$ & 240 \\ 
        $b$ & 480 \\
        \bottomrule
    \end{tabular}
    
    \caption{Propriétés du matériau}
    \label{tab:material_propertiesTAB5}
\end{table}

\begin{formal}
    \textbf{Question 2}\\
    Réponse à la question concernant le comportement du matériau...
\end{formal}

% Ajout d'une section supplémentaire pour voir l'effet dans l'index
\section{Analyse des Résultats}
Cette section est ajoutée pour démontrer le fonctionnement de la table des matières.

\end{document}
